\section{Drawing-grounded verification}
\label{sec:method}

\begin{figure}[t]
\centering
\begin{tikzpicture}[font=\scriptsize, node distance=3mm and 3mm,
  box/.style={draw, rounded corners=2pt, align=center, minimum height=7mm, fill=#1},
  box/.default={white}, >={Stealth[length=1.6mm]}]
\node[box=blue!6] (svg) {Edited SVG\\(any source)};
\node[box=orange!10, right=of svg] (geo) {Visible geometry\\and text only};
\node[box=orange!10, right=of geo] (net) {Recovered\\network};
\node[box=green!8, below=of net] (solve) {MNA / Newton\\$+$ Hardy Cross};
\node[box=green!8, left=of solve] (check) {Criteria from\\drawing notes};
\node[box=gray!10, left=of check] (score) {Edit score:\\same network?};
\draw[->] (svg) -- (geo);
\draw[->] (geo) -- node[above]{\tiny nodes} (net);
\draw[->] (net) -- (solve);
\draw[->] (solve) -- (check);
\draw[->] (check) -- (score);
\node[box=blue!6, above=2mm of svg] (editor) {Editor: source SVG $+$ instruction $\to$ patch};
\draw[->] (editor) -- (svg);
\end{tikzpicture}
\caption{\textbf{\ours.} Any edited drawing is read as built: geometry and labels become a network, the
network is solved, and the result is checked against the criteria printed on the drawing and compared with
the intended design. Identifiers and metadata are never read.}
\label{fig:pipeline}
\end{figure}

\paragraph{Task.} A source drawing $S$ of a network and an instruction $u$ (\eg ``Add R7 = 10\,$\Omega$ in
parallel with R1'') define a target design $D^\star$. An editor returns a drawing $\hat S$. We ask two
questions of $\hat S$: does it describe $D^\star$ (\emph{edit correctness}), and does the network it describes
satisfy the criteria stated on the drawing (\emph{verdict})? Both are decided by recovering a network
$\mathcal{R}(\hat S)$ from the drawing and solving it.

\subsection{Recovering a network from a drawing}
Recovery parses the SVG with transforms applied and rejects scripts, external references, clipping and
masking. It then reads only what is visibly drawn.

\paragraph{Connectivity.} Stroked straight segments are candidate conductors (for pipe networks, only heavy
strokes). A segment endpoint lying in the interior of another segment is a T-junction and splits it; a filled
junction dot splits every segment passing through its centre. Two segments that merely cross stay
unconnected, which is the drafting convention for both schematics and piping plans. Endpoints closer than
$\tau = 10^{-3}$ drawing units are merged. For circuits, all points joined by wire form one electrical node.
For pipe networks, nodes are points of degree $\neq 2$, dotted points and points on a reservoir outline, and
every chain of segments between two nodes is one pipe. A $90^\circ$ turn along a chain is an elbow; other
bend angles are rejected.

\paragraph{Components.} A rectangle whose two short-side midpoints coincide with conductor endpoints is a
resistor between the corresponding nodes. A circle is classified by where conductors meet it, not by its size:
if a conductor passes through its centre it is a junction dot, and if exactly two conductors end on its rim it
is a voltage source, whose positive terminal is the one nearer the ``$+$'' mark inside it. Filled rectangles
touched by pipe ends are reservoirs. Classifying by incidence rather than radius or stroke width keeps recovery
unchanged under uniform scaling of the drawing.

\paragraph{Values and criteria.} Labels are parsed with fixed grammars (\eg \code{R3 4.7 k$\Omega$},
\code{P5 320 m \O 150}, \code{J2 q=6 L/s z=10 m}) and assigned one-to-one to the nearest component or pipe,
within a reach proportional to the drawing's own scale (the median resistor body or pipe segment). Design
criteria are read from the drawing's notes: the resistor power rating and supply fuse; the pipe roughness,
elbow loss coefficient, minimum pressure head and maximum velocity. A drawing that leaves any component
unlabelled, or has two labels competing for one component, is reported as unreadable rather than guessed at.

\subsection{Solving the recovered network}
\paragraph{Circuits.} We use modified nodal analysis~\cite{ho1975mna} with one grounded node per connected
part of the network, so a floating fragment carries no current and does not make the system singular. A source
whose terminals share a node, or a loop made only of sources, makes the system singular; such a drawing is
reported as ill-posed. The outputs are every branch current, every resistor's dissipation and every source's
current; the criteria are $P_R \le P_\text{rated}$ for each resistor and $|I_V| \le I_\text{fuse}$ for each
source.

\paragraph{Pipe networks.} For pipe $p$ from node $a$ to node $b$ with flow $Q_p$, the head loss is
\begin{equation}
h_p(Q_p) = \Big(f(\mathrm{Re}_p,\varepsilon/D_p)\,\frac{L_p}{D_p} + n_p K_\text{elbow}\Big)
\frac{v_p |v_p|}{2g},
\end{equation}
where $v_p = 4Q_p/(\pi D_p^2)$, $f$ is Churchill's friction factor~\cite{churchill1977} and $n_p$ is the number
of elbows counted from the drawing. We solve jointly for pipe flows and junction heads,
\begin{equation}
h_p(Q_p) = H_a - H_b \;\;\forall p, \qquad
\textstyle\sum_{\text{in}} Q_p - \sum_{\text{out}} Q_p = q_j \;\;\forall j,
\end{equation}
with reservoir heads fixed, by Newton's method with backtracking~\cite{todini1988}, stopping at residuals below
$10^{-10}$\,m and $10^{-13}$\,m$^3$/s. Because Churchill's $f$ tends to $64/\mathrm{Re}$ in laminar flow, $h_p$
has a finite, non-zero slope at $Q_p = 0$ and the Jacobian stays regular when a flow reverses. As an independent
check, every network is also solved by Hardy Cross loop corrections~\cite{hardycross1936}, with multiple
reservoirs joined through a virtual datum node. The criteria are $|v_p| \le v_\text{max}$ for each pipe and
$H_j - z_j \ge p_\text{min}$ for each junction.

\subsection{Scoring an edited drawing}
An edit is \emph{correct} when $\mathcal{R}(\hat S)$ equals the target network $D^\star$ in two senses. The
first is topological: we compare, for every node, the multiset of named component terminals it touches, which
is a complete invariant of a network with named components, together with every component value and every
criterion. The second is physical: the two solutions agree, with currents within $10^{-9}$ relative plus
1\,pA, flows within $10^{-9}$\,m$^3$/s and heads within $10^{-7}$\,m. Markup is irrelevant: a drawing that
re-orders elements, renames identifiers or re-formats numbers is credited if it describes the same network,
and a drawing that is byte-for-byte plausible but disconnects a wire by $0.01$ units is not.

\subsection{Validating the verifier}
\label{sec:validation}
\Cref{tab:controls} lists the analytical controls. They include a balanced and an unbalanced Wheatstone
bridge; we initially chose bridge values that happened to balance, which made the ``unbalanced'' control
vacuous and hid a sign error, and we record this as a caution. Beyond the controls, \roundTripTotal{}
independently generated drawings (source designs and all their edits, over two seeds) recover to exactly
the network they were drawn from, as do all \datasetDrawings{} drawings in the released datasets.
Unit tests confirm four further properties: recovery is unchanged by stripping identifiers, shuffling
elements, translating and uniformly scaling the drawing by $\times0.25$ to $\times4$; a T-junction joins,
while a crossing joins only when dotted; moving only the ``$+$'' mark reverses every branch current; and moving
one wire endpoint by $0.01$ units breaks the connection. Two of these tests found real defects in the first
version. A fixed-radius rule for junction dots misclassified dots once the drawing was scaled up, and an
absolute current tolerance misjudged an edit whose two sources exactly cancel.

\begin{table}[t]
\centering
\footnotesize
\setlength{\tabcolsep}{3pt}
\begin{tabular}{@{}p{4.55cm}rr@{}}
\toprule
Control & \ours & Reference \\
\midrule
\controlRows
\bottomrule
\end{tabular}
\caption{\textbf{Analytical controls} for the two solvers. References are closed forms, an independent
bisection, or a second algorithm.}
\label{tab:controls}
\end{table}
